\documentclass[11pt,a4paper]{article}

\usepackage{kotex}
\usepackage{fontspec}
\setmainfont{Times New Roman}
\setmainhangulfont{AppleMyungjo}
\setsanshangulfont{Apple SD Gothic Neo}

\usepackage[a4paper,top=18mm,bottom=18mm,left=20mm,right=20mm]{geometry}
\usepackage{fancyhdr}
\setlength{\headheight}{24pt}
\usepackage{enumitem}
\usepackage{mdframed}
\usepackage{array}
\usepackage{booktabs}

\setlength{\parindent}{1em}
\linespread{1.18}

\pagestyle{fancy}
\fancyhf{}
\renewcommand{\headrulewidth}{0.6pt}
\fancyhead[L]{\sffamily\bfseries 2026 고1 영어 변형 — 지문 19}
\fancyhead[R]{\sffamily Daniel \& 기차 지연}
\renewcommand{\footrulewidth}{0pt}
\fancyfoot[C]{\framebox[16mm][c]{\thepage}}

\newcommand{\qno}[1]{\par\vspace{8pt}\noindent{\sffamily\bfseries #1.}\ }
\newcommand{\instr}[1]{{\sffamily #1}\par\vspace{3pt}}
\newlist{choices}{enumerate}{1}
\setlist[choices]{label=,leftmargin=1.5em,itemsep=2pt,topsep=3pt,parsep=0pt}
\newmdenv[linewidth=0.6pt,roundcorner=3pt,innertopmargin=7pt,innerbottommargin=7pt,
          innerleftmargin=9pt,innerrightmargin=9pt,skipabove=5pt,skipbelow=5pt]{pbox}
\newcommand{\nemo}[1]{\fbox{\,#1\,}}

\begin{document}

\begin{center}
{\fontsize{15}{17}\selectfont\sffamily\bfseries [지문 19] 다음 글을 읽고 물음에 답하시오.}
\end{center}
\vspace{4pt}

\begin{pbox}
Daniel was on his \underline{way} to the train station for a business trip.
He was riding a bus when it suddenly stopped.
The driver announced, ``We have a \underline{mechanical} problem.''
Daniel grew \rule{2.5cm}{0.4pt} about missing his train.
Without \underline{hesitation}, he got off the bus and took a taxi to the station.
Still, he kept \underline{worrying} about not making it in time.
When he finally reached the station, he found out the train was delayed by ten minutes.
Once he settled into his seat, \underline{he let out a long breath of relief}.
\end{pbox}

\vspace{6pt}

\qno{1}\instr{다음 빈칸에 들어갈 말로 가장 적절한 것은?}
\begin{choices}
\item ① indifferent \hfill ② apprehensive \hfill ③ elated
\item ④ relieved \hfill ⑤ triumphant
\end{choices}

\vspace{6pt}

\qno{2}\instr{밑줄 친 `he let out a long breath of relief'가 윗글에서 의미하는 바로 가장 적절한 것은?}
\begin{choices}
\item ① he was exhausted after running to the station
\item ② he felt proud of his quick decision to take a taxi
\item ③ he felt relaxed after learning the train had not left without him
\item ④ he was disappointed that the train was not on time
\item ⑤ he was relieved because the mechanical problem was fixed
\end{choices}

\vspace{6pt}

\qno{3}\instr{[서술형] 다음 문장을 우리말로 자연스럽게 해석하시오.}

\vspace{4pt}
\begin{pbox}
Without hesitation, he got off the bus and took a taxi to the station.
\end{pbox}

\vspace{4pt}
\noindent $\Rightarrow$ \rule{\linewidth}{0.4pt}

% ===================== Q4 =====================
\qno{4}\instr{Daniel의 심경 변화로 가장 적절한 것은?}
\begin{choices}
\item ① frustrated $\rightarrow$ indifferent
\item ② anxious $\rightarrow$ relieved
\item ③ confident $\rightarrow$ disappointed
\item ④ embarrassed $\rightarrow$ proud
\item ⑤ excited $\rightarrow$ regretful
\end{choices}

\end{document}
